% =============================================================
% Paper – „Herleitungen im historischen Kontext“
% Autor: Christian Weilharter
% Serie: Theoretische Physik – Herleitungen
% Ziel: 6×9 in (KDP/Print) + EPUB-kompatibles PDF (Apple Books)
% =============================================================

\documentclass[12pt,oneside]{scrartcl}

% ---- Sprache & Kodierung ----
\usepackage[utf8]{inputenc}
\usepackage[T1]{fontenc}
\usepackage[ngerman]{babel}
% ---- Literatur (BibLaTeX/Biber) ----
\usepackage[backend=biber,style=numeric,sorting=nyt]{biblatex}
\renewcommand*{\bibfont}{\small}
\addbibresource{../bib/literatur.bib} 
% ---- Gemeinsamer Stil ----
\usepackage{../styles/paper-style-de}

% =============================================================
% Serien- und Publikationsmetadaten
% =============================================================
\newcommand{\SeriesTitle}{Theoretische Physik – Herleitungen im historischen Kontext}
\newcommand{\PaperTitle}{Der Compton-Effekt – Herleitung, Geschichte und Anwendungen}
\newcommand{\PaperAuthor}{Christian Weilharter}
\newcommand{\PaperDate}{26.\ Oktober\ 2025}

\newcommand{\PaperDOI}{\emph{DOI (Zenodo)}: \texttt{10.5281/zenodo.xxxxxxx}}
\newcommand{\PaperISBNPDF}{\emph{ISBN (Print)}: 978-3-xxxx-xxxx-x}
\newcommand{\PaperISBNEPUB}{\emph{ISBN (E-Book)}: 978-3-xxxx-xxxx-x}
\newcommand{\PaperKeywords}{Compton-Effekt; Quantenphysik; Röntgenstreuung; Photonenimpuls; Vierervektoren}

% ---------- Metadatenblock ----------
\newcommand{\PaperMetadata}{
	\vspace{1em}
	\begin{flushleft}
		\small
		\setlength{\parskip}{3pt}
		\noindent\PaperDOI\\
		\PaperISBNPDF\\
		\PaperISBNEPUB\\
		\textit{Schlagwörter:} \PaperKeywords
	\end{flushleft}
	\vspace{1em}
}

% =============================================================
% Titel
% =============================================================
\title{\PaperTitle\\[0.5em]\large \SeriesTitle}
\author{\PaperAuthor}
\date{\PaperDate}

% =============================================================
% Dokument
% =============================================================
\begin{document}
	
	\thispagestyle{empty}
	\begin{center}
		{\Huge\bfseries \PaperTitle\par}
		\vspace{0.5em}
		{\Large\itshape \SeriesTitle\par}
		\vspace{2em}
		{\Large \PaperAuthor\par}
		\vspace{0.5em}
		{\large \PaperDate\par}
	\end{center}
	
	\PaperMetadata
	
	\vspace{1em}
	\hrule
	\vspace{1em}
	
	\noindent\textbf{Kurzfassung.}  
	Dieses Paper verbindet eine kompakte Herleitung der Compton-Formel mit dem historischen Kontext von Planck bis Compton und skizziert Anwendungen in der modernen Physik. Der Fokus liegt auf Verständlichkeit, mathematischer Präzision und der Einbettung in die Wissenschaftsgeschichte.
	
	\vspace{1em}
	\hrule
	\vspace{2em}
	
	% =============================================================
	% Inhalt
	% =============================================================
	
	\section{Einleitung}
	
	Als Arthur Compton im Jahr 1923 die Streuung von Röntgenstrahlung an Elektronen untersuchte, stieß er auf ein Resultat, das die klassische Wellentheorie des Lichts nicht erklären konnte. Die Wellenlänge der gestreuten Strahlung war größer als die der einfallenden – und die Differenz hing vom Streuwinkel ab. Dieses Experiment lieferte den ersten direkten Nachweis für den Impulsübertrag zwischen Photon und Elektron und bestätigte damit den Teilchencharakter des Lichts.
	
	Der Compton-Effekt markiert einen Wendepunkt in der Physik: Er verbindet Energie- und Impulserhaltung auf quantenmechanischer Ebene und führt zur bekannten \textbf{Compton-Formel}, die den Wellenlängen-Shift präzise beschreibt. Die folgende Darstellung zeigt diese Herleitung Schritt für Schritt – historisch eingeordnet, mathematisch fundiert und physikalisch nachvollziehbar.
	
	\vspace{1.2em}
	\begin{quote}
		\small
		\textbf{Hinweis zur Methodik.}  
		Dieses Paper wurde in enger Zusammenarbeit zwischen Mensch und Künstlicher Intelligenz erstellt. 
		Die physikalische Argumentation und Struktur stammen vollständig vom Autor; die KI diente als Werkzeug zur Formulierung, Präzisierung und didaktischen Gestaltung. 
		Ziel ist es, zu zeigen, wie moderne KI-Unterstützung wissenschaftliche Arbeit nicht ersetzt, sondern bereichert – durch Klarheit, Struktur und Transparenz.
	\end{quote}
	\vspace{1.2em}
	

\section{Historischer Kontext – Der Weg zum Compton-Effekt}

\subsection{Ausgangslage: Das Licht als Welle}

Zu Beginn des 20.~Jahrhunderts galt Licht nach der Theorie von \textbf{James Clerk Maxwell} (1860er Jahre) eindeutig als elektromagnetische Welle. 
Dieses Modell konnte Interferenz, Beugung und Polarisation hervorragend erklären – und schien vollständig.

Doch ab dem Jahr \textbf{1900} begann das Weltbild zu wanken:

\begin{itemize}
	\item \textbf{Max Planck} führte das Wirkungsquantum \( h \) ein, um die Schwarzkörperstrahlung zu erklären.
	\item \textbf{Albert Einstein} (1905) interpretierte dieses \( h\nu \) als reale Energieportionen des Lichts – die „Lichtquanten“.
	\item Sein \emph{Photoeffekt} zeigte, dass Licht teilchenartige Eigenschaften haben muss.
\end{itemize}

Die Physik war plötzlich zwischen zwei widersprüchlichen Bildern gefangen: 
\emph{Welle oder Teilchen?}
\vspace{1em}
\begin{HistoryBox}[Zitat von Albert Einstein (1909) \cite{Einstein1909Radiation}]
	\begin{quote}
		\emph{"We have two contradictory pictures of reality; separately neither of them fully explains the phenomena of light, but together they do."}
	\end{quote}
	
	\hfill --- \textbf{Albert Einstein}, \emph{On the Development of Our Views Concerning the Nature and Constitution of Radiation} (1909)
	\label{box:Einstein1909}
\end{HistoryBox}



\vspace{1em}
\subsection{Die Jahre vor Compton}

Zwischen 1905 und 1920 blieb das Lichtquantenkonzept umstritten. 
Viele Physiker (darunter \textbf{Robert Millikan}) bestätigten zwar Einsteins Formeln experimentell, glaubten aber nicht an die Existenz realer Photonen.  
Man betrachtete Einsteins Modell als „nützliches Rechenmittel“, jedoch nicht als physikalische Realität.

Parallel entwickelten sich Röntgenstrahlen zur neuen Präzisionsmethode in der Atomphysik.  
Ihre Streuung an Elektronen wurde mit der klassischen Elektrodynamik (\emph{Thomson-Streuung}) erklärt – allerdings nur für langwellige Strahlung.  
Bei kurzen Wellenlängen (Röntgenbereich) traten unerklärliche Abweichungen auf.

\vspace{1em}
\subsection{Der Durchbruch: Arthur Holly Compton (1922–1923)}

Der junge amerikanische Physiker \textbf{Arthur H.~Compton} untersuchte die Streuung von Röntgenstrahlen an Graphit.  
Dabei beobachtete er ein verblüffendes Phänomen:

\begin{quote}
	Nach der Streuung änderte sich die Wellenlänge der Röntgenstrahlen – und zwar \emph{abhängig vom Streuwinkel}.
\end{quote}

Diese Verschiebung ließ sich nicht durch klassische Wellendynamik erklären.  
Compton deutete sie als \emph{Stoß zwischen einem Photon und einem Elektron}, ganz analog zu elastischer Teilchenstreuung.  
Daraus folgerte er die heute berühmte Beziehung:

\[
\lambda' - \lambda = \frac{h}{m_e c}(1 - \cos\theta)
\]

Diese Formel bestätigte erstmals quantitativ die \textbf{Teilchennatur des Lichts}.

\vspace{1em}
\subsection{Die Reaktionen der Fachwelt}

Die Veröffentlichung von \textbf{1923} schlug in der Fachwelt ein wie eine Bombe:  
Viele Physiker, die Einstein 1905 noch misstraut hatten, erkannten nun: \emph{Das Photon ist real.}  
Compton erhielt 1927 den \textbf{Nobelpreis für Physik} für seine Entdeckung.  

Der Effekt wurde bald in zahlreichen Laboren bestätigt – unter anderem durch \textbf{Duane}, \textbf{Shimizu}, \textbf{Debye} und \textbf{Waller}.  
Damit war das Lichtquant endgültig etabliert – und der Weg zur \textbf{Quantenmechanik der 1920er Jahre} 
(Heisenberg, Schrödinger, Dirac) war frei.

\vspace{1em}
\subsection{Bedeutung}

Der Compton-Effekt war weit mehr als nur ein Experiment:

\begin{itemize}
	\item Er verknüpfte klassische Mechanik (\emph{Impulserhaltung}) mit der Quantentheorie (\emph{Photonenergie}).
	\item Er bewies, dass elektromagnetische Strahlung \emph{Impuls trägt}.
	\item Er zeigte, dass die Welt der Teilchen und Wellen \emph{nicht trennbar}, sondern \emph{komplementär} ist.
\end{itemize}

Der Compton-Effekt markierte somit den entscheidenden Schritt von der klassischen zur modernen Physik – 
ein Wendepunkt, der das physikalische Weltbild grundlegend veränderte.

	
	\section{Physikalische Grundlagen}
	Photonenimpuls, Energie- und Impulserhaltung, Streugeometrie (Definition des Winkels $\theta$). 
	Kurze Herleitung von $p=\tfrac{h}{\lambda}$, Beziehung zu Energie $E=hc/\lambda$.
	
	\begin{figure}[h]
		\centering
		\fbox{\parbox{0.85\linewidth}{\centering Abbildung: Streugeometrie des Compton-Effekts (Platzhalter).}}
		\caption{Einfallendes Photon, Streuwinkel $\theta$, Austrittsrichtung des Elektrons.}
		\label{fig:geometry}
	\end{figure}
	
	\section{Herleitung der Compton-Formel}
	Energie- und Impulserhaltung vor und nach dem Stoß, Eliminieren der Elektronengeschwindigkeit, Übergang zur Wellenlängenform:
	\begin{equation}
		\Delta\lambda = \lambda' - \lambda = \frac{h}{m_e c}\,(1 - \cos\theta),
		\label{eq:compton}
	\end{equation}
	mit der Compton-Wellenlänge $\lambda_C = \tfrac{h}{m_e c}$.
	
	\section{Beispielrechnung}
	Wähle z.\,B. $\lambda = \SI{0.071}{\nano\meter}$ und $\theta=90^\circ$. 
	Setze \cref{eq:compton} ein und diskutiere die Größenordnung von $\Delta\lambda$.
	
	\section{Diskussion und Grenzen}
	Thomson-Grenzfall ($\lambda \gg \lambda_C$), relativistische Korrekturen bei hohen Energien, 
	experimentelle Bestätigungen und typische Fehlerquellen.
	
	\section{Anwendungen}
	Röntgenstreuung, Materialanalyse (z.\,B. Bestimmung der Elektronendichte), astrophysikalische Phänomene.
	
	\section{Fazit}
	Knappe Zusammenfassung (ca. 150–200 Wörter): 
	Was zeigt der Compton-Effekt über die Natur des Lichts? 
	Relevanz für moderne Physik.
	
	\section*{Danksagung (optional)}
	Kurzer Dank an Mitwirkende, Institutionen, Förderhinweise.
	
	% =============================================================
	% Literatur
	% =============================================================
	\renewcommand*{\bibfont}{\small}
	
	\printbibliography
	% =============================================================
	% Rückseite
	% =============================================================
	\clearpage
	\thispagestyle{empty}
	
	\vspace*{1.5cm}
	
	{\Large\bfseries Der Compton-Effekt}\\[0.5em]
	{\large\itshape Herleitung, Geschichte und Anwendungen}
	
	\vspace{1.2cm}
	
	\noindent
	\textbf{Über dieses Paper}\\
	Der Compton-Effekt markiert einen Wendepunkt in der Physik:
	Zum ersten Mal ließ sich eindeutig zeigen, dass Licht nicht
	nur Welleneigenschaften besitzt, sondern auch Impuls überträgt –
	wie ein Teilchen. Dieses Paper stellt die Herleitung der
	Compton-Formel verständlich und prägnant dar, ergänzt durch den
	historischen Kontext und die Bedeutung für die moderne Physik.
	
	\vspace{1.0cm}
	
	\noindent
	\textbf{Open Access}\\
	Die \emph{kostenlose} PDF-Version ist verfügbar unter:\\[2pt]
	\href{https://doi.org/10.5281/zenodo.xxxxxxx}{\texttt{https://doi.org/10.5281/zenodo.xxxxxxx}}
	
	\vspace{0.8cm}
	
	\noindent
	\textbf{Serie}\\
	\textit{Theoretische Physik – Herleitungen im historischen Kontext}\\
	Weitere Hefte und Materialien:\\
	\href{https://mathandphysics.de}{\texttt{mathandphysics.de}}
	
	\vfill
	
	\noindent
	\small
	\textit{Bibliografische Angaben}\\
	\PaperDOI\\
	\PaperISBNPDF\\
	\PaperISBNEPUB
	
	\vspace{0.8cm}
	
	\noindent
	\footnotesize
	© 2025 Christian Weilharter — Alle Rechte vorbehalten.\\
	Gedruckt auf Abruf über Amazon KDP. Satz: \LaTeX.
	
	\vfill
	\begin{center}
		\small
		\textit{Math \& Physics Paper Series – Open Access unter}\\[2pt]
		\href{https://mathandphysics.de}{\texttt{https://mathandphysics.de}} \, | \,
		\href{https://zenodo.org/communities/mathandphysics}{\texttt{Zenodo: Math \& Physics Open Science Collection}}
	\end{center}
	
\end{document}
